\section{El acoplamiento de $Q$ a la materia y el caso $\delta Q=0$}
\label{sec:acoplamiento}

En la sección~\ref{sec:numerico} se eligió acoplar $Q$ solo al inflatón, con $Q(N)=Q_i\ee^{-\gamma N}$, y la tabla~\ref{tab:decisiones} decía que era la única opción consistente con un campo. Esa afirmación merece una revisión, y hay además un punto de partida distinto en la literatura del grupo, el de \cite{leon22} y \cite{piccirilli23}, donde $Q$ es homogéneo ($\delta Q=0$). Esta sección examina qué implican ambos puntos de vista. Las comprobaciones simbólicas que la sostienen están en \texttt{symbolic\_checks.py} (cuatro funciones, cuatro pruebas automáticas); las derivaciones a mano siguen siendo \etq{P}.

\subsection{Verificación simbólica de la ecuación de los tensores}
\label{sec:simbTT}

La sección~\ref{sec:tensores} obtuvo a mano que $h_{ij}$ obedece la ecuación de RG. Se verificó con SymPy el paso central: para la métrica $g=\mathrm{diag}\big(-1,\,a^2(1+\varepsilon h),\,a^2(1-\varepsilon h),\,a^2\big)$ con $h=h(t,z)$ (una polarización, transversa y sin traza a primer orden), el tensor de Einstein cumple
\begin{equation}
G^x{}_x-G^y{}_y=\varepsilon\Big(\ddot h+3H\dot h-\frac{\partial_z^2h}{a^2}\Big)+O(\varepsilon^2),
\label{eq:GxxGyy}
\end{equation}
y las componentes espaciales fuera de la diagonal se anulan. En UG, las ecuaciones $G^a{}_b+\Lb\,\delta^a_b=\kap T^a{}_b$ tienen $\Lb$ escalar, de modo que $\Lb\,\delta^a_b$ se cancela en la diferencia $x$--$y$; si además $T^i{}_j$ no tiene parte TT (el supuesto de la sección~\ref{sec:supuestos}), \eqref{eq:GxxGyy} debe anularse, y eso es la ecuación de RG. \emph{Alcance}: se comprueba la ecuación a primer orden. No se comprueba la acción a segundo orden ni la cancelación del término del vínculo (secciones~\ref{sec:accion2} y siguientes), que sigue verificada solo a mano. El signo de \eqref{eq:GxxGyy} se determinó con el propio cálculo (mi expectativa inicial del signo era la opuesta); no cambia la ecuación.

\subsection{Un escalar canónico con $Q=Q(\phi)$: UG como RG con otro potencial}
\label{sec:unescalar}

Para un campo escalar canónico con potencial $V$, y con el tensor energía-impulso usado en este trabajo, se satisface la identidad
\begin{equation}
\nabla_aT^{ab}=(\Box\phi-V')\,\nabla^b\phi,
\label{eq:divT}
\end{equation}
que se comprobó simbólicamente a primer orden, en una ansatz escalar perturbado sin E (lapse, shift y curvatura arbitrarios, dependientes de $t$ y de una coordenada espacial), para $V$ cuadrático y de Starobinsky. En UG la difusión se define por $\nabla_aT^{ab}=\nabla^bQ$ (ecuación~\eqref{eq:Qdef}). Si la materia es un único escalar canónico, \eqref{eq:divT} obliga a que $\nabla_bQ$ sea proporcional a $\nabla_b\phi$ \emph{en todas las componentes}, no solo la temporal. Eso implica que $Q$ es función de $\phi$ \emph{localmente, donde $\nabla\phi\neq0$} (extender la conclusión globalmente requiere hipótesis adicionales):
\begin{equation}
Q=Q(\phi),\qquad \Box\phi=V'(\phi)+Q'(\phi),\qquad \delta Q=Q'(\phi)\,\delta\phi.
\label{eq:QdePhi}
\end{equation}
Aquí y en el resto de la sección $V_{\rm eff}$ se escribe en unidades $\MP=1$ (en unidades explícitas, $V_{\rm eff}=V+Q+\Lambda_0\MP^2$). En el fondo, $\dot Q=Q'\dot\phi$ y la ecuación \eqref{eq:KGQ} se convierte en $\ddot\phi+3H\dot\phi+V'+Q'=0$; junto con \eqref{eq:H2UG}, $3H^2=\tfrac12\dot\phi^2+V+Q(\phi)+\Lambda_0$, son \emph{exactamente} las ecuaciones de RG con el potencial efectivo
\begin{equation}
V_{\rm eff}(\phi)=V(\phi)+Q(\phi)+\Lambda_0 .
\label{eq:Veff}
\end{equation}
Se comprobó simbólicamente que $\dot H=-\dot\phi^2/2$ se sigue de \eqref{eq:H2UG} y \eqref{eq:KGQ}, y que con $\dot Q=Q'\dot\phi$ la ecuación de Klein--Gordon de UG coincide con la de RG con $V_{\rm eff}$. Para las perturbaciones, \eqref{eq:QdePhi} sale de la componente espacial de \eqref{eq:divT} (comprobada a primer orden: $\partial_x\delta Q=E_0\,\partial_x\delta\phi$, con $E_0=\Box\phi-V'$ de fondo), y de que las ecuaciones de campo $G^a{}_b+\Lb\delta^a_b=\kap T^a{}_b$ con $\Lb=\Lambda_0+\kap Q(\phi)$ son las de RG con $V_{\rm eff}$. Esa última equivalencia, a nivel de las perturbaciones escalares, la argumento a mano y no la verifiqué simbólicamente. Tampoco discuto el gauge restringido de UG (transformaciones que conservan $f\,\mathrm d^4x$, con $\partial_a(f\xi^a)=0$) ni la compatibilidad de las descripciones con campo uniforme con el vínculo: no se hizo el chequeo explícito.

\begin{resultado}[Consecuencia, con sus condiciones]
Para un escalar canónico y una difusión local especificada como ley $Q(\phi)$, las ecuaciones de UG pueden reescribirse como las de RG con el potencial $V_{\rm eff}=V+Q(\phi)+\Lambda_0$, en el fondo (verificado simbólicamente) y, argumentado a mano, en las perturbaciones. Las corridas de las secciones~\ref{sec:resultados} y~\ref{sec:potenciales} usan $Q(N)$; su interpretación como $Q(\phi)=Q(N(\phi))$ requiere reconstruir y fijar esa función sobre la rama monótona considerada, y entonces comparan RG con $V$ y RG con $V_{\rm eff}(\phi)=V(\phi)-Q_i\big(1-\ee^{-\gamma N(\phi)}\big)$ sobre esa trayectoria. Una función fijada como ley del modelo y una función reconstruida sobre una trayectoria no son lo mismo: la segunda no define automáticamente la misma teoría para cualquier condición inicial o perturbación fuera de esa trayectoria. No se afirma una equivalencia irrestricta de teorías, ni de su cuantización, a partir de esta cuenta.

\emph{Sobre el final de la inflación.} En el cuadrático, la inflación de UG termina antes que la de RG porque la constante negativa $\Lambda_0=-Q_i$ resta a la energía efectiva. Con $K=\dot\phi^2/2$ y $U=V_{\rm eff}$ es $3H^2=K+U$ y $\epsilon_1=3K/(K+U)$, de modo que el criterio implementado, $\epsilon_1=1$, equivale a $U=2K$, no a $U=0$: el final \emph{no coincide} con el cero de $V_{\rm eff}$ (donde $V=Q_i$, y solo si $Q$ residual es despreciable). Una versión anterior de este recuadro identificaba ambas cosas; era una imprecisión de la explicación, no del criterio numérico, que estaba bien definido.
\end{resultado}

Dos consecuencias prácticas, ambas condicionadas. Primera: para ese modelo el sector escalar \emph{no está derivado} en este trabajo; si la equivalencia perturbativa es correcta, valdrían las fórmulas estándar de RG con $V_{\rm eff}$, incluida $P_\mathcal{R}=H^2/(8\pi^2\MP^2\epsilon_1)$ a orden cero con $\epsilon_1=-\dot H/H^2$ del fondo de UG, y los valores $r_{\rm std}$ y $n_{s,{\rm std}}$ de la tabla~\ref{tab:igualN} serían las predicciones a orden cero de ese modelo, pero eso no se verificó (la equivalencia perturbativa se argumentó a mano). Segunda: en esa rama el resultado no tiene contenido propio de UG; lo que hace es reparametrizar el potencial.

\subsection{El caso $\delta Q=0$}
\label{sec:dQcero}

Las referencias \cite{leon22,piccirilli23} suponen $Q$ homogéneo, $\delta Q=0$, y una materia de radiación pura. La componente espacial de la relación anterior da, para un escalar canónico y $\delta Q=0$ impuesto en el marco especificado,
\[
E_0\,\partial_x\delta\phi=0 .
\]
\emph{Si la perturbación del inflatón es inhomogénea} ($\partial_x\delta\phi\neq0$), entonces $E_0=0$ y $\dot Q=E_0\,\dot\phi=0$: en esas condiciones un inflatón con $\delta Q=0$ solo admite $Q$ constante (la UG conservativa, equivalente a RG con $\Lambda=\Lambda_0$). El control simbólico comprueba la componente espacial y la temporal de fondo de \eqref{eq:divT} y resuelve la ecuación en esa rama genérica (con $\partial_x\delta\phi$ declarado no nulo); \emph{no} comprueba ni descarta la otra rama.

Este resultado tiene tres límites:
\begin{itemize}[leftmargin=1.6em,itemsep=2pt]
\item \emph{El producto también se anula si $\partial_x\delta\phi=0$.} Para un campo espacialmente uniforme la relación no fuerza $E_0=0$.
\item \emph{$\delta Q=0$ no es invariante de gauge.} Para escalares, bajo un cambio infinitesimal del tiempo $t\to t+\xi^0$, $\widetilde{\delta\phi}=\delta\phi-\dot{\bar\phi}\,\xi^0$ y $\widetilde{\delta Q}=\delta Q-\dot{\bar Q}\,\xi^0$ \etq{E}. Imponer $\delta Q=0$ requiere entonces una especificación de gauge o una condición física adicional. En el caso $Q=Q(\phi)$, una descripción con campo uniforme también hace uniforme a $Q$, sin que eso anule por sí solo su evolución de fondo. En UG debe examinarse, además, la compatibilidad de esa descripción con las transformaciones permitidas (difeomorfismos que conservan $f\,\mathrm d^4x$, con $\partial_a(f\xi^a)=0$) y con el vínculo; no basta trasladar una elección de gauge de RG. Esto no se hizo.
\item \emph{El alcance de la conclusión.} Vale en el marco especificado y para perturbaciones inhomogéneas del inflatón. La extensión a otros gauges o a perturbaciones de campo uniforme requiere un análisis adicional.
\end{itemize}
Por eso, lo que se afirma es: «para perturbaciones inhomogéneas del inflatón en el marco especificado, imponer $\delta Q=0$ exige $E_0=0$ y, por tanto, $\dot Q=0$». Se retira la afirmación general de que, en el marco de $\delta Q=0$, la materia \emph{no puede} ser un inflatón que intercambia energía con $Q$. Para un fluido perfecto, con $u^a u_a=-1$, $h_b{}^c=\delta_b{}^c+u_bu^c$, $D_b=h_b{}^c\nabla_c$ y $a_b=u^c\nabla_cu_b$, la proyección espacial es
\[
(\rho+p)a_b+D_bp=D_bQ.
\]
La transferencia de momento se mide en el marco del fluido. Aunque $Q$ sea homogéneo en coordenadas, a primer orden puede resultar $D_iQ=u_i\dot{\bar Q}$; por tanto no se anula en general para un fluido con velocidad perturbada. Elegir radiación como modelo estudiado sigue siendo legítimo: lo que se retira es presentarla como una necesidad demostrada por este argumento.

Implicaciones para este trabajo:
\begin{itemize}[leftmargin=1.6em,itemsep=2pt]
\item El resultado analítico tensorial (secciones~\ref{sec:tensores} y~\ref{sec:simbTT}) no depende de esta discusión: solo usa que $Q$ es escalar y que $T^i{}_j$ no tiene parte TT, y vale para radiación y para un inflatón.
\item Lo que \emph{sí} depende es el fondo. El modelo numérico de este trabajo (inflatón con $Q(N)$ prescripta) se interpreta como \eqref{eq:QdePhi} sobre la rama reconstruida (para perturbaciones no uniformes puede tener $\delta Q\neq0$ y equivale a RG con $V_{\rm eff}$ sobre esa trayectoria) o como una prescripción fenomenológica del fondo sin sector escalar. No es el fondo de \cite{leon22}, en el que, según lo leído para este trabajo, el contenido es radiación y la aceleración proviene de $Q(t)$ (sin inflatón).
\item Para estudiar la elección de materia de las referencias con $\delta Q=0$, se adoptó el fondo de radiación más $Q(t)$: $3H^2=\rho_r+Q+\Lambda_0$ con $\dot\rho_r+4H\rho_r=-\dot Q$. El código de los modos no cambia, porque solo usa $H(N)$ y $\epsilon_1(N)$; se implementaron el fondo y sus controles en la sección~\ref{sec:leon}. El sector escalar en ese marco es el de \cite{leon22,piccirilli23}, con el factor $3^{3/2}$ sin resolver (secciones~\ref{sec:leonescalar} y~\ref{sec:abiertos}).
\end{itemize}

Se decidió usar el marco de $\delta Q=0$ con radiación, el de \cite{leon22,piccirilli23} (sección~\ref{sec:leon}); el caso del inflatón con $Q=Q(\phi)$ de esta sección se conserva como resultado, con las condiciones dichas, sobre la consistencia del modelo de las secciones~\ref{sec:resultados} y~\ref{sec:potenciales}.
